\documentclass[14pt,a4paper]{extarticle}
\usepackage{fontspec}
\usepackage{polyglossia}
\setmainlanguage{russian}
\setmainfont{Times New Roman}
\setmonofont{Courier New}
\newfontfamily\cyrillicfonttt{Courier New}
\usepackage{setspace}
\onehalfspacing
\usepackage{geometry}
\geometry{left=30mm, right=15mm, top=20mm, bottom=20mm}
\usepackage{indentfirst}
\setlength{\parindent}{1.25cm}
\setlength{\parskip}{0pt}
\usepackage{xurl}
\urlstyle{same}
\usepackage{hyperref}
\usepackage{booktabs}
\usepackage{ragged2e}
\usepackage{graphicx}
\usepackage{float}
\usepackage{pdflscape}
\usepackage{listings}
\usepackage{xcolor}
\usepackage{array}
\usepackage{amsmath}
\usepackage{wasysym}
\allowdisplaybreaks[3]

\lstset{
language=Haskell,
basicstyle=\footnotesize\ttfamily,
keywordstyle=\color{blue},
stringstyle=\color{red},
commentstyle=\color{green!50!black},
showstringspaces=false,
breaklines=true,
frame=single,
numbers=left,
numberstyle=\tiny\color{gray},
stepnumber=1,
keepspaces=true,
extendedchars=true,
inputencoding=utf8
}

\newcommand{\num}[1]{#1}

\begin{document}
\sloppy
\thispagestyle{empty}
\centering{
Министерство образования и науки Российской Федерации\\
Федеральное государственное автономное образовательное учреждение высшего образования\\
«Санкт-Петербургский политехнический университет Петра Великого»\\[10pt]
Институт компьютерных наук и кибербезопасности\\
Высшая школа технологий искусственного интеллекта\\
02.03.01 Математика и компьютерные науки\\
\vspace{\fill}
{\large
Отчёт по дисциплине\\
{\Large \bfseries Вычислительная математика}\\
Лабораторная работа №5\\
{\Large <<Численное интегрирование>>}\\
Вариант 18\\
\vspace{\fill}
}
{
\flushright
{\bfseries Выполнил:}\\
студент группы 5130201/40003\\
Джабраилов А.В.\\
\hspace{9cm}Подпись: \hrulefill \\
{\bfseries Принял:}\\
К. физ.-мат. наук, доц., доц.\\
Пак В.Г.\\
\hspace{9cm}Подпись: \hrulefill \\
\hspace{9cm}<<\rule{1.5cm}{0.4pt}>> \hrulefill\ 2026~г.
}
\centering{
Санкт-Петербург, 2026
}
}
\justifying
\newpage
\sloppy

\section*{1 Задание}
\begin{enumerate}
    \item Вычислить приближённо определённый интеграл по простейшим формулам при различных шагах разбиения, оценить погрешности приближения.
    \item Вычислить приближённо значение интеграла по формулам Ньютона--Котеса с различными количествами узлов.
    \item Вычислить приближённо значение интеграла по формулам Гаусса с различными количествами узлов, оценить погрешности.
\end{enumerate}

\section*{2 Цель работы}
Освоить квадратурные формулы численного интегрирования, реализовать простейшие формулы, формулы Ньютона--Котеса и формулы Гаусса для заданного варианта, а также выполнить оценку погрешности там, где она задаётся формулами методических указаний.

\newpage

\section*{3 Ход работы}
Для варианта 18 необходимо вычислить интеграл от функции
\[
    y=\operatorname{arcctg} 2^x
\]
на отрезке
\[
    x\in[1;3].
\]
В программе используется равенство
\[
    \operatorname{arcctg} t = \arctan\frac{1}{t}, \qquad t>0.
\]
Так как $2^x>0$ на всём рассматриваемом отрезке, подынтегральная функция имеет вид
\[
    f(x)=\arctan\frac{1}{2^x}.
\]

\begin{table}[H]
\centering
\caption{Исходные данные варианта 18.}
\begin{tabular}{|c|c|}
\hline
Параметр & Значение \\ \hline
Функция & $\operatorname{arcctg}2^x$ \\ \hline
Отрезок интегрирования & $[1;3]$ \\ \hline
$n$ для формулы Гаусса & $5$ \\ \hline
$n$ для формулы Ньютона--Котеса & $4$ \\ \hline
\end{tabular}
\end{table}

Для выполнения лабораторной работы была создана программа на языке \textit{Haskell}. В программе не используются готовые численные интеграторы и не используется адаптивный метод. Все вычисления выполнены по формулам, приведённым в методических указаниях.

\subsection*{3.1 Простейшие квадратурные формулы}
Ниже приведён фрагмент кода, реализующий простейшие квадратурные формулы.

\begin{lstlisting}[caption={Фрагмент кода простейших квадратурных формул}]
leftRect n func l r =
    let h = (r - l) / fromIntegral n
        xs = [l + fromIntegral i * h | i <- [0 .. n - 1]]
    in h * sum (map func xs)

rightRect n func l r =
    let h = (r - l) / fromIntegral n
        xs = [l + fromIntegral i * h | i <- [1 .. n]]
    in h * sum (map func xs)

middleRect n func l r =
    let h = (r - l) / fromIntegral n
        xs = [l + (fromIntegral i + 0.5) * h | i <- [0 .. n - 1]]
    in h * sum (map func xs)

trapezoid n func l r =
    let h = (r - l) / fromIntegral n
        innerSum = sum [func (l + fromIntegral i * h) | i <- [1 .. n - 1]]
    in h * ((func l + func r) / 2.0 + innerSum)
\end{lstlisting}

По методическим указаниям отрезок $[a;b]$ разбивается на $n$ частичных отрезков равной длины. Шаг и узлы разбиения имеют вид
\[
    h=\frac{b-a}{n}, \qquad x_i=a+ih, \qquad i=0,1,\ldots,n.
\]

Формула левых прямоугольников:
\[
    I \approx h\sum_{i=0}^{n-1} f(x_i).
\]
Подставим $a=1$, $b=3$, $n=4$:
\[
    h=\frac{3-1}{4}=0{,}5,
    \qquad x_i=1+0{,}5i.
\]
\[
    I^{(4)}_{\text{left}}
    =0{,}5\left(f(1)+f(1{,}5)+f(2)+f(2{,}5)\right)
    =0{,}611716113624.
\]

Формула правых прямоугольников:
\[
    I \approx h\sum_{i=1}^{n} f(x_i).
\]
Подставим $a=1$, $b=3$, $n=4$:
\[
    I^{(4)}_{\text{right}}
    =0{,}5\left(f(1{,}5)+f(2)+f(2{,}5)+f(3)\right)
    =0{,}442069806397.
\]

Формула центральных прямоугольников:
\[
    I \approx h\sum_{i=0}^{n-1} f\left(x_i+\frac{h}{2}\right).
\]
Подставим $a=1$, $b=3$, $n=4$:
\[
    I^{(4)}_{\text{middle}}
    =0{,}5\left(f(1{,}25)+f(1{,}75)+f(2{,}25)+f(2{,}75)\right)
    =0{,}520882771903.
\]

Формула трапеций:
\[
    I \approx h\left(\frac{f(a)+f(b)}{2}+\sum_{i=1}^{n-1} f(x_i)\right).
\]
Подставим $a=1$, $b=3$, $n=4$:
\[
    I^{(4)}_{\text{trap}}
    =0{,}5\left(\frac{f(1)+f(3)}{2}+f(1{,}5)+f(2)+f(2{,}5)\right)
    =0{,}526892960010.
\]

Формула парабол:
\[
    I \approx \frac{h}{6}\left(f(a)+f(b)+4\sum_{i=0}^{n-1} f\left(x_i+\frac{h}{2}\right)+2\sum_{i=1}^{n-1}f(x_i)\right).
\]
В программе значение по формуле парабол вычисляется на равномерной сетке с чётным числом промежутков. Для $n=4$ получено
\[
    I^{(4)}_{\text{Simp}}=0{,}522863958380.
\]

Оценки погрешностей простейших формул взяты из методических указаний:
\[
    |R|\leq \frac{b-a}{2}hM_1,
    \qquad M_1=\max_{x\in[a;b]}|f'(x)|,
\]
\[
    |R|\leq \frac{b-a}{24}h^2M_2,
    \qquad M_2=\max_{x\in[a;b]}|f''(x)|,
\]
\[
    |R|\leq \frac{b-a}{12}h^2M_2,
    \qquad M_2=\max_{x\in[a;b]}|f''(x)|,
\]
\[
    |R|\leq \frac{b-a}{180}h^4M_4,
    \qquad M_4=\max_{x\in[a;b]}|f^{(4)}(x)|.
\]
Вычисленные значения максимумов, использованные в программе:
\[
    M_1=0{,}2772588722239781,\quad
    M_2=0{,}12011325347892068,
\]
\[
    M_4=0{,}15733720319422967,\quad
    M_6=0{,}3128065157082597.
\]

\subsection*{3.2 Формулы Ньютона--Котеса}
Ниже приведён фрагмент кода, реализующий формулы Ньютона--Котеса.

\begin{lstlisting}[caption={Фрагмент кода формул Ньютона--Котеса}]
newtonCotesData n =
    case n of
        1 -> Just (1.0 / 2.0, [1.0, 1.0], 1.0 / 12.0)
        2 -> Just (1.0 / 3.0, [1.0, 4.0, 1.0], 1.0 / 90.0)
        3 -> Just (3.0 / 8.0, [1.0, 3.0, 3.0, 1.0], 3.0 / 80.0)
        4 -> Just (2.0 / 45.0, [7.0, 32.0, 12.0, 32.0, 7.0], 8.0 / 945.0)
        5 -> Just (5.0 / 288.0, [19.0, 75.0, 50.0, 50.0, 75.0, 19.0], 275.0 / 12096.0)
        _ -> Nothing

newtonCotes n func l r =
    case newtonCotesData n of
        Just (bn, coeffs, _) ->
            let h = (r - l) / fromIntegral n
                xs = [l + fromIntegral i * h | i <- [0 .. n]]
            in bn * h * sum (zipWith (*) coeffs (map func xs))
\end{lstlisting}

Формула Ньютона--Котеса из методических указаний:
\[
    \int_a^b f(x)\,dx \approx B_nh\sum_{i=0}^{n}a_i^{(n)}f(x_i),
\]
где
\[
    x_i=a+ih,\qquad i=0,1,\ldots,n,
    \qquad h=\frac{b-a}{n}.
\]

Для варианта 18 задано $n=4$. По таблице коэффициентов:
\[
    B_4=\frac{2}{45},\qquad
    a_0^{(4)}=7,
    \quad a_1^{(4)}=32,
    \quad a_2^{(4)}=12,
    \quad a_3^{(4)}=32,
    \quad a_4^{(4)}=7.
\]
Шаг равен
\[
    h=\frac{3-1}{4}=0{,}5.
\]
Узлы имеют вид
\[
    x_0=1,\quad x_1=1{,}5,\quad x_2=2,\quad x_3=2{,}5,\quad x_4=3.
\]
Подставим числа в формулу:
\[
    I_{NC}^{(4)}=\frac{2}{45}\cdot0{,}5\left(
    7f(1)+32f(1{,}5)+12f(2)+32f(2{,}5)+7f(3)
    \right).
\]
\[
    I_{NC}^{(4)}=0{,}522878949915.
\]

Остаточный член для $n=4$ по таблице методических указаний:
\[
    r_4(h)=-\frac{8h^7}{945}f^{VI}(\xi).
\]
Для оценки предельной погрешности используем модуль остаточного члена:
\[
    |r_4(h)|\leq \frac{8h^7}{945}M_6.
\]
Подставим числа:
\[
    |r_4(h)|\leq
    \frac{8\cdot0{,}5^7}{945}\cdot0{,}3128065157082597
    =2{,}068826162092\cdot10^{-5}.
\]

\subsection*{3.3 Формулы Гаусса}
Ниже приведён фрагмент кода, реализующий формулы Гаусса.

\begin{lstlisting}[caption={Фрагмент кода формул Гаусса}]
gaussData n =
    case n of
        2 -> Just [(-0.577350, 1.0), (0.577350, 1.0)]
        3 -> Just [(-0.774597, 5.0 / 9.0), (0.0, 8.0 / 9.0), (0.774597, 5.0 / 9.0)]
        4 -> Just [(-0.861136, 0.347855), (-0.339981, 0.652145),
                   (0.339981, 0.652145), (0.861136, 0.347855)]
        5 -> Just [(-0.906180, 0.236927), (-0.538469, 0.478629),
                   (0.0, 0.568889), (0.538469, 0.478629),
                   (0.906180, 0.236927)]
        _ -> Nothing

gauss n func l r =
    case gaussData n of
        Just pairs ->
            let mid = (l + r) / 2.0
                half = (r - l) / 2.0
                oneTerm (x, weight) = weight * func (mid + half * x)
            in half * sum (map oneTerm pairs)
\end{lstlisting}

В методических указаниях таблица узлов и весов Гаусса приведена для интеграла
\[
    \int_{-1}^{1}f(x)\,dx.
\]
Если интеграл вычисляется на отрезке $[a;b]$, то узлы пересчитываются по формуле
\[
    t_i=\frac{a+b}{2}+\frac{b-a}{2}x_i.
\]
Тогда
\[
    \int_a^b f(t)\,dt \approx \frac{b-a}{2}\sum_{i=1}^{n}A_if(t_i).
\]

Для варианта 18 задано $n=5$. По таблице Гаусса:
\[
\begin{array}{c|c}
    x_i & A_i \\ \hline
    -0{,}906180 & 0{,}236927 \\
    -0{,}538469 & 0{,}478629 \\
    0 & 0{,}568889 \\
    0{,}538469 & 0{,}478629 \\
    0{,}906180 & 0{,}236927
\end{array}
\]
Для отрезка $[1;3]$ имеем
\[
    \frac{a+b}{2}=2,\qquad \frac{b-a}{2}=1.
\]
Следовательно,
\[
    t_i=2+x_i.
\]
Получим узлы на отрезке $[1;3]$:
\[
    t_1=1{,}093820,\quad
    t_2=1{,}461531,\quad
    t_3=2,
    \quad t_4=2{,}538469,
    \quad t_5=2{,}906180.
\]
Подставим числа в формулу:
\begingroup\small
\[
\begin{aligned}
I_G^{(5)}={}&1\cdot\bigl(
0{,}236927f(1{,}093820)+0{,}478629f(1{,}461531)+0{,}568889f(2)\,+{}\\
&{}+0{,}478629f(2{,}538469)+0{,}236927f(2{,}906180)
\bigr).
\end{aligned}
\]
\endgroup
\[
    I_G^{(5)}=0{,}522887989168.
\]

\newpage

\section*{4 Результаты}

\subsection*{4.1 Простейшие квадратурные формулы}

\begin{table}[H]
\centering
\caption{Результаты по формуле левых прямоугольников.}
\small
\begin{tabular}{|c|c|c|}
\hline
$n$ & Значение интеграла & Оценка погрешности \\ \hline
4 & $0{,}611716113624$ & $1{,}386294361120\cdot10^{-1}$ \\ \hline
8 & $0{,}566299442763$ & $6{,}931471805599\cdot10^{-2}$ \\ \hline
16 & $0{,}544343483804$ & $3{,}465735902800\cdot10^{-2}$ \\ \hline
32 & $0{,}533553116652$ & $1{,}732867951400\cdot10^{-2}$ \\ \hline
\end{tabular}
\end{table}

\begin{table}[H]
\centering
\caption{Результаты по формуле правых прямоугольников.}
\small
\begin{tabular}{|c|c|c|}
\hline
$n$ & Значение интеграла & Оценка погрешности \\ \hline
4 & $0{,}442069806397$ & $1{,}386294361120\cdot10^{-1}$ \\ \hline
8 & $0{,}481476289150$ & $6{,}931471805599\cdot10^{-2}$ \\ \hline
16 & $0{,}501931906998$ & $3{,}465735902800\cdot10^{-2}$ \\ \hline
32 & $0{,}512347328249$ & $1{,}732867951400\cdot10^{-2}$ \\ \hline
\end{tabular}
\end{table}

\begin{table}[H]
\centering
\caption{Результаты по формуле центральных прямоугольников.}
\small
\begin{tabular}{|c|c|c|}
\hline
$n$ & Значение интеграла & Оценка погрешности \\ \hline
4 & $0{,}520882771903$ & $2{,}502359447478\cdot10^{-3}$ \\ \hline
8 & $0{,}522387524845$ & $6{,}255898618694\cdot10^{-4}$ \\ \hline
16 & $0{,}522762749500$ & $1{,}563974654673\cdot10^{-4}$ \\ \hline
32 & $0{,}522856494709$ & $3{,}909936636684\cdot10^{-5}$ \\ \hline
\end{tabular}
\end{table}

\begin{table}[H]
\centering
\caption{Результаты по формуле трапеций.}
\small
\begin{tabular}{|c|c|c|}
\hline
$n$ & Значение интеграла & Оценка погрешности \\ \hline
4 & $0{,}526892960010$ & $5{,}004718894955\cdot10^{-3}$ \\ \hline
8 & $0{,}523887865957$ & $1{,}251179723739\cdot10^{-3}$ \\ \hline
16 & $0{,}523137695401$ & $3{,}127949309347\cdot10^{-4}$ \\ \hline
32 & $0{,}522950222451$ & $7{,}819873273367\cdot10^{-5}$ \\ \hline
\end{tabular}
\end{table}

\begin{table}[H]
\centering
\caption{Результаты по формуле парабол.}
\small
\begin{tabular}{|c|c|c|}
\hline
$n$ & Значение интеграла & Оценка погрешности \\ \hline
4 & $0{,}522863958380$ & $1{,}092619466627\cdot10^{-4}$ \\ \hline
8 & $0{,}522886167939$ & $6{,}828871666416\cdot10^{-6}$ \\ \hline
16 & $0{,}522887638549$ & $4{,}268044791510\cdot10^{-7}$ \\ \hline
32 & $0{,}522887731467$ & $2{,}667527994694\cdot10^{-8}$ \\ \hline
\end{tabular}
\end{table}

\subsection*{4.2 Формулы Ньютона--Котеса}

\begin{table}[H]
\centering
\caption{Результаты по формулам Ньютона--Котеса.}
\small
\begin{tabular}{|c|c|c|}
\hline
$n$ & Значение интеграла & Оценка погрешности \\ \hline
1 & $0{,}588002603548$ & $8{,}007550231928\cdot10^{-2}$ \\ \hline
2 & $0{,}522639085352$ & $1{,}748191146603\cdot10^{-3}$ \\ \hline
3 & $0{,}522768406589$ & $7{,}769738429345\cdot10^{-4}$ \\ \hline
4 & $0{,}522878949915$ & $2{,}068826162092\cdot10^{-5}$ \\ \hline
5 & $0{,}522882856072$ & $1{,}165162894490\cdot10^{-5}$ \\ \hline
\end{tabular}
\end{table}

Для варианта 18 по формуле Ньютона--Котеса с $n=4$ получено
\[
    I_{NC}^{(4)}=0{,}522878949915,
    \qquad |r_4(h)|\leq2{,}068826162092\cdot10^{-5}.
\]

\subsection*{4.3 Формулы Гаусса}

\begin{table}[H]
\centering
\caption{Результаты по формулам Гаусса.}
\small
\begin{tabular}{|c|c|}
\hline
Количество узлов & Значение интеграла \\ \hline
2 & $0{,}523045828497$ \\ \hline
3 & $0{,}522896341738$ \\ \hline
4 & $0{,}522887598540$ \\ \hline
5 & $0{,}522887989168$ \\ \hline
\end{tabular}
\end{table}

Для варианта 18 по формуле Гаусса с пятью узлами получено
\[
    I_G^{(5)}=0{,}522887989168.
\]
Так как в методических указаниях для таблицы узлов и весов Гаусса не приведён явный остаточный член, в программе выводятся приближённые значения при разных количествах узлов.

\newpage

\section*{5 Вывод}
В ходе лабораторной работы были реализованы формулы численного интегрирования для варианта 18. Для функции $f(x)=\operatorname{arcctg}2^x$ на отрезке $[1;3]$ были вычислены приближённые значения интеграла по простейшим квадратурным формулам, по формулам Ньютона--Котеса и по формулам Гаусса.

При увеличении числа разбиений в простейших формулах значения интеграла стабилизируются около $0{,}52288$. Формулы центральных прямоугольников, трапеций и парабол дают более точные результаты, чем формулы левых и правых прямоугольников. Для заданного варианта получены основные ответы:
\[
    I_{NC}^{(4)}=0{,}522878949915,
    \qquad
    I_G^{(5)}=0{,}522887989168.
\]

\newpage

\section*{6 Приложение}
\begin{lstlisting}[caption={Файл Main.hs}, label={lst:lab6code}]
module Main where

import Text.Printf

a :: Double
a = 1.0

b :: Double
b = 3.0

f :: Double -> Double
f x = atan (1.0 / (2.0 ** x))

m1 :: Double
m1 = 0.2772588722239781

m2 :: Double
m2 = 0.12011325347892068

m4 :: Double
m4 = 0.15733720319422967

m6 :: Double
m6 = 0.3128065157082597

leftRect :: Int -> (Double -> Double) -> Double -> Double -> Double
leftRect n func l r =
    let h = (r - l) / fromIntegral n
        xs = [l + fromIntegral i * h | i <- [0 .. n - 1]]
    in h * sum (map func xs)

rightRect :: Int -> (Double -> Double) -> Double -> Double -> Double
rightRect n func l r =
    let h = (r - l) / fromIntegral n
        xs = [l + fromIntegral i * h | i <- [1 .. n]]
    in h * sum (map func xs)

middleRect :: Int -> (Double -> Double) -> Double -> Double -> Double
middleRect n func l r =
    let h = (r - l) / fromIntegral n
        xs = [l + (fromIntegral i + 0.5) * h | i <- [0 .. n - 1]]
    in h * sum (map func xs)

trapezoid :: Int -> (Double -> Double) -> Double -> Double -> Double
trapezoid n func l r =
    let h = (r - l) / fromIntegral n
        innerSum = sum [func (l + fromIntegral i * h) | i <- [1 .. n - 1]]
    in h * ((func l + func r) / 2.0 + innerSum)

simpsonComposite :: Int -> (Double -> Double) -> Double -> Double -> Double
simpsonComposite nRaw func l r =
    let n = if even nRaw then nRaw else nRaw + 1
        h = (r - l) / fromIntegral n
        oddSum = sum [func (l + fromIntegral i * h) | i <- [1,3 .. n - 1]]
        evenSum = sum [func (l + fromIntegral i * h) | i <- [2,4 .. n - 2]]
    in h / 3.0 * (func l + func r + 4.0 * oddSum + 2.0 * evenSum)

leftRightErrorBound :: Int -> Double
leftRightErrorBound n =
    let h = (b - a) / fromIntegral n
    in (b - a) * h / 2.0 * m1

middleErrorBound :: Int -> Double
middleErrorBound n =
    let h = (b - a) / fromIntegral n
    in (b - a) * h * h / 24.0 * m2

trapezoidErrorBound :: Int -> Double
trapezoidErrorBound n =
    let h = (b - a) / fromIntegral n
    in (b - a) * h * h / 12.0 * m2

simpsonErrorBound :: Int -> Double
simpsonErrorBound nRaw =
    let n = if even nRaw then nRaw else nRaw + 1
        h = (b - a) / fromIntegral n
    in (b - a) * h ** 4.0 / 180.0 * m4

newtonCotesData :: Int -> Maybe (Double, [Double], Double)
newtonCotesData n =
    case n of
        1 -> Just (1.0 / 2.0, [1.0, 1.0], 1.0 / 12.0)
        2 -> Just (1.0 / 3.0, [1.0, 4.0, 1.0], 1.0 / 90.0)
        3 -> Just (3.0 / 8.0, [1.0, 3.0, 3.0, 1.0], 3.0 / 80.0)
        4 -> Just (2.0 / 45.0, [7.0, 32.0, 12.0, 32.0, 7.0], 8.0 / 945.0)
        5 -> Just (5.0 / 288.0, [19.0, 75.0, 50.0, 50.0, 75.0, 19.0], 275.0 / 12096.0)
        _ -> Nothing

newtonCotes :: Int -> (Double -> Double) -> Double -> Double -> Double
newtonCotes n func l r =
    case newtonCotesData n of
        Nothing -> error "Unsupported Newton-Cotes n"
        Just (bn, coeffs, _) ->
            let h = (r - l) / fromIntegral n
                xs = [l + fromIntegral i * h | i <- [0 .. n]]
                values = map func xs
            in bn * h * sum (zipWith (*) coeffs values)

newtonCotesErrorBound :: Int -> Double
newtonCotesErrorBound n =
    case newtonCotesData n of
        Nothing -> error "Unsupported Newton-Cotes n"
        Just (_, _, c) ->
            let h = (b - a) / fromIntegral n
            in case n of
                1 -> c * h ** 3.0 * m2
                2 -> c * h ** 5.0 * m4
                3 -> c * h ** 5.0 * m4
                4 -> c * h ** 7.0 * m6
                5 -> c * h ** 7.0 * m6
                _ -> error "Unsupported Newton-Cotes n"

gaussData :: Int -> Maybe [(Double, Double)]
gaussData n =
    case n of
        2 -> Just
            [ (-0.577350, 1.0)
            , ( 0.577350, 1.0)
            ]
        3 -> Just
            [ (-0.774597, 5.0 / 9.0)
            , ( 0.0,      8.0 / 9.0)
            , ( 0.774597, 5.0 / 9.0)
            ]
        4 -> Just
            [ (-0.861136, 0.347855)
            , (-0.339981, 0.652145)
            , ( 0.339981, 0.652145)
            , ( 0.861136, 0.347855)
            ]
        5 -> Just
            [ (-0.906180, 0.236927)
            , (-0.538469, 0.478629)
            , ( 0.0,      0.568889)
            , ( 0.538469, 0.478629)
            , ( 0.906180, 0.236927)
            ]
        _ -> Nothing

gauss :: Int -> (Double -> Double) -> Double -> Double -> Double
gauss n func l r =
    case gaussData n of
        Nothing -> error "Unsupported Gauss n"
        Just pairs ->
            let mid = (l + r) / 2.0
                half = (r - l) / 2.0
                oneTerm (x, weight) = weight * func (mid + half * x)
            in half * sum (map oneTerm pairs)

printLine :: IO ()
printLine = putStrLn "------------------------------------------------------------"

printSimpleResult :: String -> Int -> Double -> Double -> IO ()
printSimpleResult name n value bound = do
    printf "%-20s n = %-4d value = %.12f  error_bound = %.12e\n"
        name n value bound

printNCResult :: Int -> Double -> Double -> IO ()
printNCResult n value bound = do
    printf "Newton-Cotes         n = %-4d value = %.12f  error_bound = %.12e\n"
        n value bound

printGaussResult :: Int -> Double -> IO ()
printGaussResult n value = do
    printf "Gauss                n = %-4d value = %.12f\n"
        n value

main :: IO ()
main = do
    putStrLn "Lab 5. Numerical integration"
    putStrLn "Variant 18"
    putStrLn "Function: y = arccot(2^x)"
    putStrLn "Interval: [1; 3]"
    printLine

    putStrLn "Task 1. Simple composite formulas with different partition steps"
    putStrLn "Columns: formula, n, approximate value, theoretical error bound"
    printLine

    let parts = [4, 8, 16, 32]

    mapM_ (\n -> do
            let value = leftRect n f a b
            let bound = leftRightErrorBound n
            printSimpleResult "Left rectangles" n value bound) parts

    printLine

    mapM_ (\n -> do
            let value = rightRect n f a b
            let bound = leftRightErrorBound n
            printSimpleResult "Right rectangles" n value bound) parts

    printLine

    mapM_ (\n -> do
            let value = middleRect n f a b
            let bound = middleErrorBound n
            printSimpleResult "Middle rectangles" n value bound) parts

    printLine

    mapM_ (\n -> do
            let value = trapezoid n f a b
            let bound = trapezoidErrorBound n
            printSimpleResult "Trapezoid" n value bound) parts

    printLine

    mapM_ (\n -> do
            let value = simpsonComposite n f a b
            let bound = simpsonErrorBound n
            printSimpleResult "Simpson" n value bound) parts

    printLine

    putStrLn "Task 2. Newton-Cotes formulas with different n"
    putStrLn "Columns: n, approximate value, theoretical error bound"
    printLine

    mapM_ (\n -> do
            let value = newtonCotes n f a b
            let bound = newtonCotesErrorBound n
            printNCResult n value bound) [1 .. 5]

    printLine

    putStrLn "Task 3. Gauss formulas with different number of nodes"
    putStrLn "Columns: nodes count, approximate value"
    printLine

    mapM_ (\n -> do
            let value = gauss n f a b
            printGaussResult n value) [2 .. 5]

    printLine

    putStrLn "Variant-specific results"
    let variantNC = newtonCotes 4 f a b
    let variantNCBound = newtonCotesErrorBound 4
    let variantGauss = gauss 5 f a b

    printf "Newton-Cotes n = 4 value = %.12f\n" variantNC
    printf "Newton-Cotes n = 4 error_bound = %.12e\n" variantNCBound
    printf "Gauss n = 5 value = %.12f\n" variantGauss

    printLine
    putStrLn "Done."
\end{lstlisting}

\end{document}
